\cotefile{85}
\scanpage{16}
% L01–02
donc $(\mathbb D_\infty\times\mathbb Z_2)/N\simeq\mathbb D_\infty/N\cap\mathbb D_\infty$ et on applique encore le $=$ corollaire.

% L03–10
\uline{NB} Les groupes $\deleted{\mathbb D_{\unsure{i|\infty}},}\ \mathbb D_i\times\mathbb Z_2$ ($i$ infini, ou fini pair) et $\mathbb D_i$ sont-ils effectivement plongeables \unsure{sur|dans} un $H=N(T)$ ou $B$ ? \inserted{$k$ fixé (alg. clos de \unsure{car})}
% slanted
\margin{explicitons le cas des représentations \unsure{irr. diédrales|irr. diédriques|non diédrales}}
\deleted{Pour les $\mathbb D_i$, c'est évident que oui pour $N(T)$, \unsure{prouver que ce sont alg. clos|puisque ce sont alg. clos|pourvu que le seul alg. clos} les \unsure{ss groupes|gp. des}. le cas pour $\mathbb D_i\times\mathbb Z_2$ car déjà \unsure{cas} groupe de $\mathbb Z/2\mathbb Z\times\mathbb Z/2\mathbb Z$ vient pas} \inserted{\unsure{V$_4$|$\simeq V_4$}}

% L11–13
\deleted{\uline{Proposition} \inserted{Pour $k$ alg. clos, on envisage la \unsure{c.}} Pour que $\mathbb D_i$ \unsure{admette|admettent} une repr. \unsure{projective} dans $\mathrm{GP}(1,k)$, il faut que}

% L14–16
\uline{Proposition} Soit $k$ corps alg. clos (\deleted{de} car $p$), $i\in\mathbb N^{\deleted{\unsure{*}}}$, considérons la condition
\[
(\ast)\quad\exists\ \text{repr. fidèle}\quad\mathbb D_i\lhook\joinrel\longrightarrow\mathrm{GP}(1,k)
\]

% L17–20
\uline{Alors} :

a) Si $p=0$, \inserted{ou $i\deleted{\unsure{\geq 1}}=1$}, $(\ast)$ est satisfaite. \deleted{\unsure{Si $p\ne0$ et $i=1$, $(\ast)$ est satisfaite}.}

b) \deleted{Si $p\ne0$, (i.e. $\mathbb D_i=\mathbb D_0=\mathbb D_\infty$)} \inserted{Supposons $p>0$. Si $i=0$ (i.e. $\mathbb D_i=\mathbb D_0=\mathbb D_\infty$)} $(\ast)$ équivaut à la condition : \deleted{\unsure{$k$ n'est le c.}, \illegible{1}} $k$ non algébrique sur le corps premier (ou encore : $k^*$ contient un élément d'ordre infini)

% L21–25
\deleted{\unsure{b)|c)} Si $p\ne0$, et}

c) Si $i\geq2$, $(\ast)$ équivaut à la condition :

\deleted{On a}

$(\ast a)$ On a \deleted{$i$ premier à} $p\nmid i$, ou bien $i=p$, \deleted{\unsure{et le groupe} \illegible{1}}

\deleted{ou bien $i=2p$ et $p$ impair.}

% L26–32
Dém. a), b) immédiats, en \deleted{\unsure{$(\ast)$|cas}} que $i=0$ \unsure{$(\ast)$|car}
% slanted
\margin{\unsure{explicitons pour|explicitons un peu|explicitez ceux}}
équivaut à l'existence d'un él. d'ordre infini dans $k^*$ \{ou dans $k$.\}

\deleted{Pour $i\geq2$, on note que les}

\deleted{c) On sait que les sous-groupes \inserted{d'ordre fini} de $\mathrm{GP}(1)$ sont \unsure{ou les} les groupes \unsure{cycliques et premiers à $p$, ou les}}

\deleted{b) Si \unsure{$H(k)$} \unsure{$\simeq k$ ou}, \unsure{prenant} $H=\unsure{\mathrm V_{1/2}\,k|\mathrm V_{1,2}\,k|\gamma\underset{1/2}{\cdot}U}$ $(H=\unsure{\mathrm G_m}\text{ ou }\mathrm G_a)$}
% archive sheet number: 15
